\hypertarget{objetivos}{%
\section{Objetivos}\label{objetivos}}

\begin{frame}{Objetivos}
\protect\hypertarget{objetivos-1}{}
\begin{itemize}
\item
  Classificar as otimiza\c{c}\~{o}es de código intermediário por escopo: local
  (intra-bloco), global (inter-bloco) e de la\c{c}o.
\item
  Definir formalmente e implementar dobramento de constantes, propaga\c{c}\~{a}o
  de constantes, elimina\c{c}\~{a}o de subexpress\~{o}es comuns (CSE) e elimina\c{c}\~{a}o
  de código morto (DCE).
\item
  Apresentar otimiza\c{c}\~{o}es estruturais de la\c{c}o: movimenta\c{c}\~{a}o de código
  invariante (LICM), desenrolamento e fus\~{a}o.
\item
  Definir formalmente e implementar inlining de fun\c{c}\~{o}es e elimina\c{c}\~{a}o de
  chamadas em cauda (TCO), com suas condi\c{c}\~{o}es de elegibilidade e
  corre\c{c}\~{a}o.
\end{itemize}
\end{frame}

\hypertarget{introduuxe7uxe3o}{%
\section{Introdu\c{c}\~{a}o}\label{introduuxe7uxe3o}}

\begin{frame}{O que s\~{a}o Otimiza\c{c}\~{o}es?}
\protect\hypertarget{o-que-suxe3o-otimizauxe7uxf5es}{}
\begin{itemize}
\tightlist
\item
  ``Otimiza\c{c}\~{a}o'' é um exagero: raramente o código é \textbf{ótimo}
  segundo alguma métrica global
\item
  O que se busca: código \textbf{melhor} --- mais rápido, menor, mais
  econ\^{o}mico
\item
  Aplicadas sobre a \textbf{IRT} sem alterar o significado observável do
  programa
\end{itemize}

\textbf{Classifica\c{c}\~{a}o por escopo:}

\begin{longtable}[]{@{}
  >{\raggedright\arraybackslash}p{(\columnwidth - 4\tabcolsep) * \real{0.3333}}
  >{\raggedright\arraybackslash}p{(\columnwidth - 4\tabcolsep) * \real{0.3333}}
  >{\raggedright\arraybackslash}p{(\columnwidth - 4\tabcolsep) * \real{0.3333}}@{}}
\toprule\noalign{}
\begin{minipage}[b]{\linewidth}\raggedright
Escopo
\end{minipage} & \begin{minipage}[b]{\linewidth}\raggedright
Descri\c{c}\~{a}o
\end{minipage} & \begin{minipage}[b]{\linewidth}\raggedright
Exemplos
\end{minipage} \\
\midrule\noalign{}
\endhead
\textbf{Local} (intra-bloco) & Dentro de um bloco básico & Const
folding, CSE \\
\textbf{Global} (inter-bloco) & Análise de fluxo de controle & Const
propagation global \\
\textbf{De la\c{c}o} & Exploram estrutura iterativa & LICM, unrolling,
fusion \\
\bottomrule\noalign{}
\end{longtable}
\end{frame}

\begin{frame}{Dez Otimiza\c{c}\~{o}es Neste Capítulo}
\protect\hypertarget{dez-otimizauxe7uxf5es-neste-capuxedtulo}{}
\begin{enumerate}
\tightlist
\item
  \textbf{Dobramento de constantes} (\emph{constant folding})
\item
  \textbf{Propaga\c{c}\~{a}o de constantes} (\emph{constant propagation})
\item
  \textbf{Dobramento e propaga\c{c}\~{a}o combinados}
\item
  \textbf{Elimina\c{c}\~{a}o de subexpress\~{o}es comuns} (CSE)
\item
  \textbf{Elimina\c{c}\~{a}o de código morto} (DCE)
\item
  \textbf{Movimenta\c{c}\~{a}o de código invariante de la\c{c}o} (LICM)
\item
  \textbf{Desenrolamento de la\c{c}o} (\emph{loop unrolling})
\item
  \textbf{Fus\~{a}o de la\c{c}os} (\emph{loop fusion})
\item
  \textbf{Inlining de fun\c{c}\~{o}es} (\emph{function inlining})
\item
  \textbf{Elimina\c{c}\~{a}o de chamadas em cauda} (\emph{tail-call
  elimination})
\end{enumerate}
\end{frame}

\begin{frame}{Terminologia}
\protect\hypertarget{terminologia}{}
Uma express\~{a}o é \textbf{pura} se n\~{a}o produz efeitos colaterais:
\[e \text{ pura} \iff e \in \{\mathbf{CONST}, \mathbf{TEMP}, \mathbf{NAME}\} \text{ ou } e = \mathbf{BINOP}(\mathit{op}, e_1, e_2) \text{ com } e_1, e_2 \text{ puras}\]

\begin{itemize}
\tightlist
\item
  \(\mathbf{MEM}\), \(\mathbf{CALL}\), \(\mathbf{ESEQ}\) \textbf{n\~{a}o s\~{a}o
  puras}
\item
  \(\mathit{freeTemps}(e)\) = temporários que ocorrem livres em \(e\)
\item
  \(\mathit{defs}(s)\) = temporários \textbf{escritos} por \(s\)
\item
  \(\mathit{uses}(s)\) = temporários \textbf{lidos} por \(s\)
\end{itemize}
\end{frame}

\hypertarget{dobramento-de-constantes}{%
\section{Dobramento de Constantes}\label{dobramento-de-constantes}}

\begin{frame}{Defini\c{c}\~{a}o Formal}
\protect\hypertarget{definiuxe7uxe3o-formal}{}
\[\mathit{fold}(\mathbf{BINOP}(\mathit{op}, e_1, e_2)) = \begin{cases}
  \mathbf{CONST}(\mathit{op}(n_1, n_2)) & \text{se } e_1' = \mathbf{CONST}(n_1) \text{ e } e_2' = \mathbf{CONST}(n_2) \\
  \mathit{simplify}(\mathit{op}, e_1', e_2') & \text{caso contrário}
\end{cases}\]

\textbf{Identidades algébricas} aplicadas por \(\mathit{simplify}\):

\[e + 0 = e \qquad e - 0 = e \qquad e \times 0 = 0 \qquad e \times 1 = e\]

\textbf{Cuidado}: divis\~{a}o por zero em compile-time deve ser evitada!
\end{frame}

\begin{frame}[fragile]{Exemplo de Dobramento}
\protect\hypertarget{exemplo-de-dobramento}{}
Antes:

\begin{lstlisting}
BINOP(ADD, CONST(3), BINOP(MUL, CONST(2), CONST(5)))
\end{lstlisting}

Depois:

\begin{lstlisting}
CONST(13)
\end{lstlisting}

Também:

\begin{lstlisting}
BINOP(MUL, TEMP(x), CONST(0))  →  CONST(0)
\end{lstlisting}

(sem precisar saber o valor de \passthrough{\lstinline!x!})

\textbf{Corre\c{c}\~{a}o}:
\(\sigma \vdash e \Downarrow v \iff \sigma \vdash \mathit{fold}(e) \Downarrow v\)
\end{frame}

\hypertarget{propagauxe7uxe3o-de-constantes}{%
\section{Propaga\c{c}\~{a}o de
Constantes}\label{propagauxe7uxe3o-de-constantes}}

\begin{frame}{Ambiente de Propaga\c{c}\~{a}o}
\protect\hypertarget{ambiente-de-propagauxe7uxe3o}{}
Um \textbf{ambiente de propaga\c{c}\~{a}o}
\(\rho : \mathit{Temp} \rightharpoonup \mathbb{Z}\) mapeia temporários a
seus valores constantes conhecidos.

\[\mathit{foldProp}_\rho(\mathbf{TEMP}(t)) = \begin{cases}
  \mathbf{CONST}(\rho(t)) & \text{se } t \in \mathit{dom}(\rho) \\
  \mathbf{TEMP}(t) & \text{caso contrário}
\end{cases}\]

Para comandos --- ao processar \(\mathbf{MOVE}(\mathbf{TEMP}(t), e)\):

\[\rho' = \begin{cases}
  \rho[t \mapsto n] & \text{se } e' = \mathbf{CONST}(n) \\
  \rho \setminus \{t\} & \text{caso contrário}
\end{cases}\]
\end{frame}

\begin{frame}[fragile]{Exemplo de Propaga\c{c}\~{a}o Combinada}
\protect\hypertarget{exemplo-de-propagauxe7uxe3o-combinada}{}
Antes:

\begin{lstlisting}
MOVE(TEMP(x), CONST(5))
MOVE(TEMP(y), BINOP(MUL, TEMP(x), CONST(2)))
RETURN(TEMP(y))
\end{lstlisting}

Depois (com \(\rho = \{x \mapsto 5, y \mapsto 10\}\)):

\begin{lstlisting}
MOVE(TEMP(x), CONST(5))
MOVE(TEMP(y), CONST(10))
RETURN(CONST(10))
\end{lstlisting}

Combinado com DCE, as duas primeiras linhas s\~{a}o eliminadas:
\passthrough{\lstinline!RETURN(CONST(10))!}
\end{frame}

\hypertarget{eliminauxe7uxe3o-de-subexpressuxf5es-comuns-cse}{%
\section{Elimina\c{c}\~{a}o de Subexpress\~{o}es Comuns
(CSE)}\label{eliminauxe7uxe3o-de-subexpressuxf5es-comuns-cse}}

\begin{frame}{Defini\c{c}\~{a}o}
\protect\hypertarget{definiuxe7uxe3o}{}
A \textbf{tabela CSE}
\(\theta : \mathit{PureExpr} \rightharpoonup \mathit{Temp}\) mapeia
express\~{o}es puras ao temporário que armazena seu valor mais recente.

Para \(\mathbf{MOVE}(\mathbf{TEMP}(t), e)\) onde \(e\) é pura:

\[\mathit{CSE}_\theta(\mathbf{MOVE}(\mathbf{TEMP}(t), e)) = \begin{cases}
  \mathbf{MOVE}(\mathbf{TEMP}(t), \mathbf{TEMP}(\theta(e))) & \text{se } e \in \mathit{dom}(\theta) \\
  \mathbf{MOVE}(\mathbf{TEMP}(c), e);\; \mathbf{MOVE}(\mathbf{TEMP}(t), \mathbf{TEMP}(c)) & \text{caso contrário}
\end{cases}\]

A tabela é \textbf{reiniciada} em fronteiras de bloco (LABEL, JUMP,
CJUMP, RETURN).
\end{frame}

\begin{frame}[fragile]{Exemplo de CSE}
\protect\hypertarget{exemplo-de-cse}{}
Antes:

\begin{lstlisting}
MOVE(TEMP(a), BINOP(ADD, TEMP(x), TEMP(y)))
MOVE(TEMP(b), BINOP(ADD, TEMP(x), TEMP(y)))
\end{lstlisting}

Depois:

\begin{lstlisting}
MOVE(TEMP(_cse0), BINOP(ADD, TEMP(x), TEMP(y)))
MOVE(TEMP(a),     TEMP(_cse0))
MOVE(TEMP(b),     TEMP(_cse0))
\end{lstlisting}

A express\~{a}o \(x + y\) é calculada \textbf{uma única vez}.

\textbf{Por que só express\~{o}es puras?} \passthrough{\lstinline!MEM(e)!}
pode mudar após escrita em memória.
\end{frame}

\hypertarget{eliminauxe7uxe3o-de-cuxf3digo-morto-dce}{%
\section{Elimina\c{c}\~{a}o de Código Morto
(DCE)}\label{eliminauxe7uxe3o-de-cuxf3digo-morto-dce}}

\begin{frame}[fragile]{Fase 1: Código Inalcan\c{c}ável}
\protect\hypertarget{fase-1-cuxf3digo-inalcanuxe7uxe1vel}{}
Uma instru\c{c}\~{a}o é \textbf{inalcan\c{c}ável} se está imediatamente após JUMP ou
RETURN sem LABEL intermediário:

\begin{lstlisting}
RETURN(CONST(42))
MOVE(TEMP(x), CONST(1))   ← inalcan\c{c}ável! (sem LABEL antes)
LABEL(continue)
\end{lstlisting}
\end{frame}

\begin{frame}{Fase 2: Atribui\c{c}\~{o}es Mortas}
\protect\hypertarget{fase-2-atribuiuxe7uxf5es-mortas}{}
Uma instru\c{c}\~{a}o \(\mathbf{MOVE}(\mathbf{TEMP}(t), e)\) é \textbf{morta} se
\(t\) é redefinido antes de ser lido, ou a fun\c{c}\~{a}o termina sem usar
\(t\).

Estado de análise: \((\mathit{pend}, \mathit{dead})\)

\begin{itemize}
\tightlist
\item
  \(\mathit{pend}\) = atribui\c{c}\~{o}es pendentes (ainda n\~{a}o consumidas)
\item
  \(\mathit{dead}\) = índices das instru\c{c}\~{o}es marcadas para remo\c{c}\~{a}o
\end{itemize}

Em fronteiras de bloco (LABEL, JUMP, CJUMP): descarte conservador de
\(\mathit{pend}\).
\end{frame}

\begin{frame}[fragile]{Exemplo de DCE}
\protect\hypertarget{exemplo-de-dce}{}
Após propaga\c{c}\~{a}o de constantes:

\begin{lstlisting}
MOVE(TEMP(x), CONST(5))     ← x nunca lido → morta
MOVE(TEMP(y), CONST(10))    ← y nunca lido → morta
RETURN(CONST(10))
\end{lstlisting}

Depois da DCE:

\begin{lstlisting}
RETURN(CONST(10))
\end{lstlisting}
\end{frame}

\hypertarget{otimizauxe7uxf5es-de-lauxe7o}{%
\section{Otimiza\c{c}\~{o}es de La\c{c}o}\label{otimizauxe7uxf5es-de-lauxe7o}}

\begin{frame}[fragile]{Estrutura de La\c{c}o na IRT}
\protect\hypertarget{estrutura-de-lauxe7o-na-irt}{}
Um la\c{c}o \passthrough{\lstinline!while!} na IRT linearizada tem o padr\~{a}o:

\[\underbrace{\mathbf{LABEL}(\ell_h)}_{\text{cabe\c{c}a}} \;\; \mathbf{CJUMP}(e, \ell_b, \ell_f) \;\; \underbrace{\mathbf{LABEL}(\ell_b)}_{\text{corpo}} \;\; s_1 \cdots s_k \;\; \mathbf{JUMP}(\mathbf{NAME}(\ell_h)) \;\; \underbrace{\mathbf{LABEL}(\ell_f)}_{\text{saída}}\]

Um \textbf{la\c{c}o} \(L = (\ell_h, \ell_b, \ell_f, e_\mathit{cond}, B)\)
onde \(B = [s_1, \ldots, s_k]\).
\end{frame}

\begin{frame}{Movimenta\c{c}\~{a}o de Invariante de La\c{c}o (LICM)}
\protect\hypertarget{movimentauxe7uxe3o-de-invariante-de-lauxe7o-licm}{}
Uma instru\c{c}\~{a}o \(\mathbf{MOVE}(\mathbf{TEMP}(t), e)\) é
\textbf{invariante de la\c{c}o} se:

\begin{enumerate}
\tightlist
\item
  \(e\) é pura
\item
  \(\mathit{freeTemps}(e) \cap \mathit{defs}(L) = \emptyset\) ---
  nenhuma variável de \(e\) é modificada no la\c{c}o
\item
  \(t\) é definido exatamente \textbf{uma vez} no corpo
\item
  Nenhuma instru\c{c}\~{a}o anterior l\^{e} \(t\) no corpo
\end{enumerate}

\textbf{Transforma\c{c}\~{a}o}: instru\c{c}\~{o}es invariantes s\~{a}o i\c{c}adas
(\emph{hoisted}) para antes do cabe\c{c}alho.
\end{frame}

\begin{frame}[fragile]{Exemplo de LICM}
\protect\hypertarget{exemplo-de-licm}{}
Antes:

\begin{lstlisting}
LABEL(loop_head)
CJUMP(...)
LABEL(loop_body)
  MOVE(TEMP(k), BINOP(MUL, CONST(4), CONST(8)))  ← invariante!
  MOVE(TEMP(a), BINOP(MUL, TEMP(i), TEMP(k)))
  MOVE(TEMP(i), BINOP(SUB, TEMP(i), CONST(1)))
JUMP(NAME(loop_head))
\end{lstlisting}

Depois:

\begin{lstlisting}
MOVE(TEMP(k), BINOP(MUL, CONST(4), CONST(8)))   ← i\c{c}ado
LABEL(loop_head)
  ...
  MOVE(TEMP(a), BINOP(MUL, TEMP(i), TEMP(k)))
  MOVE(TEMP(i), BINOP(SUB, TEMP(i), CONST(1)))
\end{lstlisting}
\end{frame}

\begin{frame}[fragile]{Desenrolamento de La\c{c}o}
\protect\hypertarget{desenrolamento-de-lauxe7o}{}
Com fator \(k = 2\), replica o corpo \(k-1\) vezes extras:

\begin{lstlisting}
LABEL(head)  CJUMP(cond, body, exit)
LABEL(body)
  <corpo original>
  CJUMP(cond, body2, exit)   ← verifica\c{c}\~{a}o extra
LABEL(body2)
  <cópia do corpo>           ← renomeando rótulos internos
JUMP(NAME(head))
LABEL(exit)
\end{lstlisting}

\textbf{Benefícios}: menos avalia\c{c}\~{o}es da condi\c{c}\~{a}o, mais oportunidades
para outras otimiza\c{c}\~{o}es.
\end{frame}

\begin{frame}[fragile]{Fus\~{a}o de La\c{c}os}
\protect\hypertarget{fusuxe3o-de-lauxe7os}{}
\textbf{Loop fusion} combina dois la\c{c}os adjacentes com a \textbf{mesma
condi\c{c}\~{a}o} em um único la\c{c}o:

Antes:

\begin{lstlisting}
while (cond) { corpo1 }
while (cond) { corpo2 }
\end{lstlisting}

Depois:

\begin{lstlisting}
while (cond) { corpo1; corpo2 }
\end{lstlisting}

\textbf{Benefícios}: reduz overhead de controle, melhora localidade de
dados.

\textbf{Restri\c{c}\~{a}o}: os dois la\c{c}os devem ter a mesma condi\c{c}\~{a}o e n\~{a}o haver
depend\^{e}ncias conflitantes entre \passthrough{\lstinline!corpo1!} e
\passthrough{\lstinline!corpo2!}.
\end{frame}

\hypertarget{inlining-de-funuxe7uxf5es}{%
\section{Inlining de Fun\c{c}\~{o}es}\label{inlining-de-funuxe7uxf5es}}

\begin{frame}[fragile]{O que é Inlining?}
\protect\hypertarget{o-que-uxe9-inlining}{}
Substituir uma chamada \(\mathbf{CALL}(\mathbf{NAME}(f), \bar{a})\) pelo
corpo de \(f\):

\[\mathit{inline}(f, \bar{a}) = \mathbf{ESEQ}\!\Bigl(\mathbf{MOVE}(x_1^s, a_1) \mathbin{;} \cdots \mathbin{;} B_f^s \mathbin{;} \mathbf{LABEL}(\ell_\mathit{exit}^s),\; \mathbf{TEMP}(r^s)\Bigr)\]

\begin{itemize}
\tightlist
\item
  \(x_i^s\): par\^{a}metros \textbf{renomeados} com sufixo único \(s\)
\item
  \(B_f^s\): corpo com temporários/rótulos locais \textbf{renomeados}
  com \(s\)
\item
  Cada \passthrough{\lstinline!RETURN(e)!} em \(B_f^s\) →
  \passthrough{\lstinline!MOVE(r\^s, e); JUMP(exit\^s)!}
\end{itemize}

\textbf{Benefícios}: elimina overhead de chamada + exp\~{o}e código a
dobramento, CSE, LICM.
\end{frame}

\begin{frame}[fragile]{Elegibilidade para Inlining}
\protect\hypertarget{elegibilidade-para-inlining}{}
Uma chamada é elegível se:

\begin{enumerate}
\tightlist
\item
  \textbf{Alvo estático}: \passthrough{\lstinline!NAME(f)!}, n\~{a}o um
  ponteiro computado
\item
  \textbf{N\~{a}o auto-recursiva}: \(f \notin \mathit{calls}(B_f)\) (evita
  expans\~{a}o infinita)
\item
  \textbf{Tamanho controlado}: \(|\mathit{linearize}(B_f)| \leq \tau\)
  (evita \emph{code bloat})
\end{enumerate}

\textbf{Corre\c{c}\~{a}o}:
\(\sigma \vdash \mathbf{CALL}(\mathbf{NAME}(f), \bar{a}) \Downarrow v \iff \sigma \vdash \mathit{inline}(f, \bar{a}) \Downarrow v\)

Segue diretamente da regra sem\^{a}ntica de invoca\c{c}\~{a}o.
\end{frame}

\begin{frame}[fragile]{Implementa\c{c}\~{a}o em Haskell}
\protect\hypertarget{implementauxe7uxe3o-em-haskell}{}
\begin{lstlisting}[language=Haskell]
inlineProgram :: InlineConfig -> Program -> Program
inlineProgram cfg prog =
  removeDeadFuncs $
    map (inlineFuncDef funcMap inlineable) prog
  where
    funcMap    = Map.fromList [(funcName fd, fd) | fd <- prog]
    inlineable = Set.fromList
      [ funcName fd
      | fd <- prog
      , funcSize fd <= inlineThreshold cfg
      , not (isSelfRecursive fd) ]
\end{lstlisting}

\begin{itemize}
\tightlist
\item
  \passthrough{\lstinline!inlineThreshold!}: limiar \(\tau\) (padr\~{a}o: 10
  instru\c{c}\~{o}es)
\item
  \passthrough{\lstinline!removeDeadFuncs!}: elimina fun\c{c}\~{o}es sem
  chamadores após inlining
\item
  M\^{o}nada \passthrough{\lstinline!State Int!} gera sufixos únicos para
  cada sítio de chamada
\end{itemize}
\end{frame}

\begin{frame}[fragile]{Renomea\c{c}\~{a}o e Transforma\c{c}\~{a}o de Retorno}
\protect\hypertarget{renomeauxe7uxe3o-e-transformauxe7uxe3o-de-retorno}{}
\textbf{renameLocals}: adiciona sufixo \(s\) a todos os temporários e
rótulos locais

\textbf{transformReturns}: substitui
\passthrough{\lstinline!RETURN([e])!} por
\passthrough{\lstinline!MOVE(r\^s, e); JUMP(exit\^s)!}

\begin{lstlisting}[language=Haskell]
transformReturns retTemp exitLbl = go
  where
    go (RETURN [e]) = SEQ (MOVE (TEMP retTemp) e)
                          (JUMP (NAME exitLbl))
    go (RETURN _)   = JUMP (NAME exitLbl)
    go (SEQ s1 s2)  = SEQ (go s1) (go s2)
    go s            = s
\end{lstlisting}

Sem renomea\c{c}\~{a}o, inlinar a mesma fun\c{c}\~{a}o em dois sítios produziria
\textbf{colis\~{o}es de temporários}.
\end{frame}

\hypertarget{eliminauxe7uxe3o-de-chamadas-em-cauda}{%
\section{Elimina\c{c}\~{a}o de Chamadas em
Cauda}\label{eliminauxe7uxe3o-de-chamadas-em-cauda}}

\begin{frame}[fragile]{Chamadas em Posi\c{c}\~{a}o de Cauda}
\protect\hypertarget{chamadas-em-posiuxe7uxe3o-de-cauda}{}
Uma chamada está em \textbf{posi\c{c}\~{a}o de cauda} se o único efeito
posterior é retornar seu resultado:

\begin{longtable}[]{@{}ll@{}}
\toprule\noalign{}
Padr\~{a}o & Descri\c{c}\~{a}o \\
\midrule\noalign{}
\endhead
\passthrough{\lstinline!RETURN [CALL(NAME(f), args)]!} & Retorno com
valor \\
\passthrough{\lstinline!EXP(CALL(NAME(f), args)); RETURN []!} & Retorno
vazio \\
\bottomrule\noalign{}
\end{longtable}

\textbf{Ideia}: quando \(f = g\) (auto-recurs\~{a}o), o frame corrente pode
ser \textbf{reutilizado} --- basta atualizar os par\^{a}metros e saltar para
o início.

Resultado: recurs\~{a}o \(O(n)\) em pilha → la\c{c}o \(O(1)\) em pilha.
\end{frame}

\begin{frame}{A Transforma\c{c}\~{a}o TCO}
\protect\hypertarget{a-transformauxe7uxe3o-tco}{}
Para fun\c{c}\~{a}o \(g\) com par\^{a}metros \(\bar{x} = (x_1,\ldots,x_n)\),
introduz \(\ell_h = \mathtt{L\_tco\_}g\):

\begin{enumerate}
\tightlist
\item
  Inserir \(\mathbf{LABEL}(\ell_h)\) no início do corpo
\item
  Substituir cada chamada em cauda por:
\end{enumerate}

\[\underbrace{\mathbf{MOVE}(t_1, a_1) \mathbin{;} \cdots \mathbin{;} \mathbf{MOVE}(t_n, a_n)}_{\text{passo 1: avaliar para auxiliares}} \mathbin{;} \underbrace{\mathbf{MOVE}(x_1, t_1) \mathbin{;} \cdots \mathbin{;} \mathbf{MOVE}(x_n, t_n)}_{\text{passo 2: atualizar par\^{a}metros}} \mathbin{;} \mathbf{JUMP}(\ell_h)\]

\textbf{Por que dois passos?} Se \(\bar{a} = (x_2, x_1)\), uma única
passagem sobrescreveria \(x_1\) antes de ser lido para \(x_2\). Os
temporários \(t_i\) quebram a depend\^{e}ncia cíclica.
\end{frame}

\begin{frame}[fragile]{Exemplo: Fatorial com Acumulador}
\protect\hypertarget{exemplo-fatorial-com-acumulador}{}
Antes da TCO:

\begin{lstlisting}
RETURN [CALL(NAME(factAcc),
        [BINOP(BSub, TEMP(n),   CONST(1)),
         BINOP(BMul, TEMP(acc), TEMP(n))])]
\end{lstlisting}

Após a TCO:

\begin{lstlisting}
LABEL(L_tco_factAcc)
  ...
  MOVE(TEMP(_tco_0), BINOP(BSub, TEMP(n),   CONST(1)))
  MOVE(TEMP(_tco_1), BINOP(BMul, TEMP(acc), TEMP(n)))
  MOVE(TEMP(n),   TEMP(_tco_0))
  MOVE(TEMP(acc), TEMP(_tco_1))
  JUMP(NAME(L_tco_factAcc))
\end{lstlisting}

Recurs\~{a}o convertida em la\c{c}o --- \textbf{sem crescimento de pilha}.
\end{frame}

\begin{frame}[fragile]{Implementa\c{c}\~{a}o em Haskell}
\protect\hypertarget{implementauxe7uxe3o-em-haskell-1}{}
\begin{lstlisting}[language=Haskell]
tcoFuncDef :: FuncDef -> FuncDef
tcoFuncDef fd
  | not (anyTailCall (funcName fd) stmts) = fd
  | otherwise =
      fd { funcBody = rebuildSeq (LABEL lh : transformed) }
  where
    stmts       = linearize (funcBody fd)
    lh          = "L_tco_" ++ funcName fd
    transformed = transformStmts (funcName fd) (funcParams fd) lh stmts

makeTCO :: [Temp] -> [Expr] -> [Stmt]
makeTCO params args =
  zipWith (\t e -> MOVE (TEMP t) e) tcoTemps args ++
  zipWith (\p t -> MOVE (TEMP p) (TEMP t)) params tcoTemps ++
  [JUMP (NAME lh)]
  where tcoTemps = ["_tco_" ++ show i | i <- [0..]]
\end{lstlisting}

O pipeline completo aplica TCO \textbf{antes} do inlining, pois o
inlining pode expor novas chamadas em cauda.
\end{frame}

\hypertarget{conclusuxe3o}{%
\section{Conclus\~{a}o}\label{conclusuxe3o}}

\begin{frame}{Sumário do Capítulo}
\protect\hypertarget{sumuxe1rio-do-capuxedtulo}{}
\begin{itemize}
\tightlist
\item
  \textbf{Dobramento} e \textbf{propaga\c{c}\~{a}o} de constantes: avaliar em
  compile-time
\item
  \textbf{CSE}: evitar recomputo de express\~{o}es puras já calculadas
\item
  \textbf{DCE}: remover código inalcan\c{c}ável e atribui\c{c}\~{o}es nunca lidas
\item
  \textbf{LICM}: i\c{c}ar código invariante para fora do la\c{c}o
\item
  \textbf{Unrolling / Fusion}: otimiza\c{c}\~{o}es estruturais de la\c{c}o
\item
  \textbf{Inlining}: substituir chamada pelo corpo --- elimina overhead,
  exp\~{o}e otimiza\c{c}\~{o}es
\item
  \textbf{TCO}: converter auto-recurs\~{a}o em cauda em la\c{c}o --- pilha
  \(O(1)\)
\end{itemize}
\end{frame}

\begin{frame}{Próximos Passos}
\protect\hypertarget{pruxf3ximos-passos}{}
\begin{itemize}
\tightlist
\item
  \textbf{Gera\c{c}\~{a}o de código WebAssembly} (próximo capítulo): IRT → WAT
  executável
\item
  \textbf{Análise de fluxo de dados}: otimiza\c{c}\~{o}es globais mais precisas
\item
  \textbf{Otimiza\c{c}\~{o}es baseadas em SSA}: Single Static Assignment, uma
  defini\c{c}\~{a}o por variável
\end{itemize}
\end{frame}
